\chapter{What Exactly Is Capitalism?}
\section{An Eighteen-Yuan Milk Tea}
Pick up something you may encounter every day: a milk tea, together with its receipt.

The receipt says 18 yuan. Now do something unusual: break those eighteen yuan apart and see whose pockets they eventually enter.

The brand takes a share for franchise fees, management fees, and centrally supplied ingredients: the price of the signboard. The landlord takes another; a few square meters in a mall basement may command astonishing rent. The platform takes its commission if you ordered on your phone. Suppliers receive payment for milk, tea, tapioca pearls, cups, and straws. The employee receives a share for the minutes spent mixing your drink, paid by the hour. What remains goes to the shop owner, who may never have visited today.

Watch that last share. The employee's claim is clear: she contributed time and labor. The landlord's is reasonably clear too: he supplied premises. But what entitles the owner to his share? He mixed no tea, operated no till, perhaps never appeared. He receives money because the shop is \textbf{his}: he bought the equipment, paid the franchise fee, and bears the risk of loss.

Now watch the employee's share. Why hourly pay? Without premises, equipment, or capital, the only thing she can exchange is her time. And whether she mixes a hundred drinks or ten in an hour, her pay is unchanged. The money from those extra dozens has nothing to do with her.

There need be no villain in this drink. The owner is not necessarily unscrupulous; nobody ropes the employee to her workplace. Prices are explicit and agreement voluntary. Yet tracing all eighteen yuan from origin to destination makes questions appear on their own:

What is profit: fair compensation for supplying capital and bearing risk, or a slice taken from someone else's labor? Are voluntary transactions necessarily fair? If refusing a contract means going hungry, is agreement truly voluntary? Why do some people live by selling time and others by owning things?

Together these questions form the title's large word: \textbf{capitalism}. It may be among the words you hear most frequently in the news and find hardest to define. Some call it an economic system, others an ideology; some use it as an insult and others as praise. We will take it apart and find out what it is.

\section{Manchester at Dawn, 1842}
The best way to understand something is to return to its first large-scale appearance. Come to Manchester, England, in 1842. The following composite reconstruction draws on parliamentary investigations, medical records, and eyewitness descriptions of the time; its details can be checked.

Half past five, before dawn. Coal smoke from every household mixes with moisture, pressing yellow air over rooftops and leaving an acidic smell in the nose. A factory whistle roars like an enormous animal outside the city. Thousands of footsteps on stone streets head in one direction.

Among them is Mary, a sixteen-year-old spinner. She works twelve to fourteen hours daily, interrupted only by a few dozen minutes for meals. Fine white cotton fibers fill the workshop's shafts of light and enter workers' lungs; many develop relentless coughs before middle age. Machines are so loud that conversation requires shouting into ears, a habit workers carry home. Machinery does not wait: joining a broken thread half a second too slowly may let a racing belt take a finger. Mary knows more than one coworker missing fingers.

Factory gates lock on time. A quarter-hour's lateness costs half a day's pay. Windows rarely open, not to save glass but because warm damp air keeps cotton yarn from breaking. The environment serves machines first, people second.

Ten years earlier, Mary's father wove at home in the countryside. Work was hard, but he decided when to begin or pause for tea. Now he tends machines at another mill. He cannot say which life is better: bad harvests meant hunger in the country; in town at least wages arrive each Saturday. But he repeats one thought: in a factory, people serve machines.

A few kilometers away, mill owner Mr. Robert breakfasts in his villa. He owns three spinning mills; the newest runs a steam engine day and night. His ledger lists wages beside coal, cotton, and depreciation as just another cost. The sweat of Mary and her companions appears only as a number.

Remember this dawn. Most arguments later in the chapter grew out of this workshop.

\section{How Should This Account Be Calculated?}
Before choosing sides, lay out the conflict.

Nobody in this Manchester morning violates the law of the time. Mary comes voluntarily, having signed no contract of bondage, and can resign whenever she wants, though that means hunger. Mr. Robert pays every agreed penny. Selling inherited property and borrowing from a bank financed his mills; he bought the cotton; every penny of cloth-sale profit legally belongs to him.

Why, then, did so many workers, doctors, clergy, journalists, and even aristocrats strongly feel something was \textbf{wrong}?

Three different accusations identify the wrong. Distinguish them first; the rest of the chapter keeps returning to them.

One concerns \textbf{treatment}: excessive hours, dangerous workshops, cruel child labor. This is relatively easy to answer with shorter hours, guards, and bans on child labor. Britain subsequently did address most such problems through successive legislation.

Another concerns \textbf{distribution}: even with improved treatment, why do those weaving cloth always receive the smallest share while the biggest recipient never touches a machine? Mary earns the same for mixing a hundred teas or ten: where does the extra value go? Protective guards cannot answer a challenge to profit itself.

The third concerns \textbf{relationships}: why are all one person's means of livelihood---land, tools, premises---held by others, making her voluntary sale of time scarcely a choice? Mary's freedom is to choose freely between the mill and hunger. What kind of freedom is that?

These are capitalism's three questions: how it treats people, distributes money, and arranges relationships. Laws can repair treatment; distribution cuts deeper; relationships reach the system's foundation.

Do not answer too quickly. Before judgment, clarify capitalism itself. Does it mean markets, bosses, or increased greed? Praising or condemning an undefined word fires blanks.

\section[Mill Owners, Weavers, and Distant Looms]{The Mill Owner, the Weaver, and Looms Eight Thousand Kilometers Away}
Listen before judging. Here we must do something often neglected: \textbf{state every side's full case}, including the side you instinctively dislike.

\textbf{The first voice: the mill owner.}

Give Mr. Robert a fair hearing by developing his strongest argument, not a straw man.

First comes risk. Three mills require astronomical capital: buildings, engines, spindles, cotton inventories. Cloth prices fluctuate annually. A depression, shipping disaster, or competitor's newer machinery could wipe him out. Before profit reaches his pocket, the chance of loss is real. His share first compensates the possibility of losing everything.

Second come organization and waiting. Coordinating hundreds of workers, tonnes of materials, and complex machines into an orderly production line is genuine, scarce labor. Moreover, money invested in cotton may take a year or two to return after spinning, shipping, and sale. Ordinary people cannot or will not wait; he will.

Third is the hardest argument: without mills, Mary would not even have this job. In 1835, British scholar Andrew Ure published The Philosophy of Manufactures, systematically defending factories. Broadly, he argued that factories gave the poor steadier incomes and better lives than the countryside and that machinery relieved rather than increased toil. Much of it fails scrutiny today, but one fact deserves recognition: rural poverty before factories was also harsh and often less secure.

This is no weak case. Profit is not spontaneously appearing stolen goods; it corresponds to real functions: risk bearing, organization, and advance funding while waiting. Any alternative eliminating profit must explain who performs them.

\textbf{The second voice: the weaver.}

Mary's case is equally complete. First, time: life consists of time, and surrendering fourteen hours daily to machines surrenders being human itself, leaving only eating and sleeping to maintain production. Second, autonomy: her grandfather could stop when tired or tend a sick child; now the whistle replaces her will, and even a breathing pause requires permission from machinery. Third, contribution: workers' hands weave every inch, while the owner's touch no thread. Why does his name come first in the ledger?

These voices went beyond complaint. Between 1811 and 1816, English Midlands weavers launched a famous movement, entering mills at night and destroying machinery. Their legendary leader Ned Ludd gave them the name Luddites. Later mockery depicts ignorant enemies of progress, but many were skilled workers. They attacked not machinery as an object but the arrangement behind it: machines lowering wages, children replacing skilled labor, and people becoming appendages. They smashed visible machines because they could not reach the invisible system.

\textbf{The third voice: eight thousand kilometers away.}

Manchester's machines made cloth too cheap for handweaving to compete. Ships carried it worldwide, notably to India, one of cotton textiles' oldest centers. For millennia Indian fine cotton reached Europe and western Asia; Dhaka muslin could pass through a ring. Nineteenth-century machine cloth flooded in, silencing handlooms and destroying innumerable livelihoods. India changed from cloth exporter to importer.

Notice this voice's distinctive position. Mary at least entered a mill voluntarily; Indian weavers never even had that opportunity. Nobody contracted with them. Their livelihood was shattered by a force they had never seen and could not negotiate with. Understanding capitalism requires hearing more than the contract's signatories.

Together the voices produce an uncomfortable picture: each tells part of the truth. Robert's risk, Mary's lost time, and Indian weavers' desperation are all real. This is not evasive everybody-has-a-point compromise. It identifies a system producing efficiency and costs simultaneously, distributed to different people.

\section{Thinking Bridge: Five Components of Capitalism}
After all this argument, what does capitalism mean? Build a portable thinking tool: \textbf{unpacking a word}. Before loving or hating a large term, separate its components, inspect each, then reassemble them.

The capitalist machine has at least five components.

\textbf{Component one: markets.} Markets are places and mechanisms of voluntary exchange: the school-gate shop, bargaining over vegetables, or selling used textbooks online. Their simple rule is that both parties must see advantage for a transaction to occur.

\textbf{Component two: capital.} Capital is not synonymous with money. Holiday gift money in a piggy bank is money but not capital: lying there, it does not grow. Wealth becomes capital when invested in production to generate more wealth. Household savings are not capital; using them to buy premises, stock, and labor in pursuit of more money turns them into capital. Its defining character is \textbf{the pursuit of expansion}.

\textbf{Component three: wage labor.} A person freely sells working time for wages and has no other livelihood---no land, premises, or capital. Mary is a wage laborer. Both free and propertyless matter: unfree means slave or serf; other means of livelihood would let her choose not to go.

\textbf{Component four: profit.} Revenue remaining after all costs belongs to capital's owner. You saw it in the milk-tea receipt: compensation for risk and organization, and the disputed surplus. Who should receive it?

\textbf{Component five: the enterprise.} This organizes the other four: a legally recognized person, the company, owning capital, employing labor, selling products, and distributing profits to owners. A tea franchise, spinning mill, and short-video company are all enterprises.

The five components are assembled. But real thought begins now: \textbf{having components is not enough; ask which dominate}.

The essential counterexample is easily missed: \textbf{markets, merchants, and profit do not automatically mean capitalism}.

Markets are much older. Over two thousand years ago, when Sima Qian wrote Records of the Grand Historian, Chang'an and Luoyang markets bustled and Silk Road caravans traveled. Over a millennium ago, Arab and Persian merchants crossed the Indian Ocean, exchanging spices, porcelain, and cloth among three continents and devising partnership-like contracts. Late-medieval Venice and Genoa accumulated vast merchant fortunes; spices enriched Venice's whole city-state. Ming and Qing Jiangnan had dense market towns and elaborate silk and porcelain networks. Edo Osaka's Dojima rice market traded rice bills in what is often called an early futures market. These worlds had markets, merchants, profits, and even banks and negotiable instruments.

Yet historians generally do not call them capitalist societies. Why? Their \textbf{organization of production} was not dominated by wage labor. Owner-farmers and tenants cultivated, households wove, and guild masters worked with apprentices. Profit chiefly came from trading margins rather than organizing hired production. Markets were society's marketplace, not its engine.

Capitalism's novelty was assembling the five into a main engine: \textbf{production itself is chiefly organized by capital owners hiring free, propertyless workers for profit, which becomes new capital in an expanding cycle}. This engine first took shape around early-modern Europe and accelerated fully in the Industrial Revolution. Manchester's whistle sounded its full-speed operation.

Test your new word-unpacking tool:

\begin{itemize}
\item Is Northern Song Bianjing's morning market, with shouting vendors and fluttering tavern banners, capitalism? (It has markets, but household and guild production: no.)
\item A middle-school student earns three hundred yuan selling chilled jelly from a holiday stall. Is he a capitalist? (He works himself without employees: no, he is self-employed.)
\item What if he works the stall but also hires two classmates hourly and keeps the proceeds? (The components begin to assemble: tiny in scale, but the structure is there.)
\end{itemize}
Unpacking lets you hear more than an emotional blur in capitalism. Ask which component is meant, how dominant it is, and who receives its efficiencies and costs.

\section[A Chinese Saying and Scottish Books]{Evidence: One Chinese Saying and Two Scottish Books}
Read some primary material. On why people pursue gain and what should restrain it, Chinese writers over two millennia ago and a Scotsman over two centuries ago left checkable written words.

First China. Western Han historian Sima Qian wrote in the Biographies of Money-makers within Records of the Grand Historian:

\begin{quote}
All the world's bustling comes for gain; all its jostling goes for gain.
\end{quote}
People hurry to and fro for benefit. Notice his attitude: he is \textbf{describing}, not scolding. The same chapter discusses governing the desire for wealth: best follow this nature, next guide it through incentives, then teach, then restrain through law; worst compete with the people for gain. Two millennia ago, someone already treated self-interest as a natural force to accommodate rather than preach away, asking how institutions should arrange it.

Now Scotland. In 1776, Adam Smith published An Inquiry into the Nature and Causes of the Wealth of Nations, usually The Wealth of Nations. A frequently quoted passage, as conveyed in standard Chinese translations, reads:

\begin{quote}
We expect dinner not from the kindness of the butcher, brewer, or baker, but from their concern for their own advantage. We appeal not to their benevolence but to their self-interest, speaking not of our needs but of their benefit.
\end{quote}
Isolated, it is often offered as proof that Smith preferred greater selfishness. Reading only this turns a whole person into a cartoon. Three additions are necessary.

First, Smith taught \textbf{moral philosophy} at Glasgow, not economics, which did not yet exist as a discipline. Seventeen years before The Wealth of Nations came The Theory of Moral Sentiments. Its opening, broadly in standard Chinese translations, says however selfish people are thought, their nature contains principles making them care about others' fortunes and need others' happiness, even when their only reward is pleasure at seeing it. Smith called this \textbf{sympathy}: not pity but imagining oneself in another's circumstances, the foundation of morality.

Second, Smith distrusted merchants more deeply than textbooks suggest. A famous passage in Book I says people of the same trade rarely meet even for amusement without ending in conspiracy against the public or a plan to raise prices. He devoted much space to attacking monopolies, exclusionary guild rules, and merchant lobbying for government favors.

Third, markets were not his answer to everything. The final book identifies tasks markets handle poorly and public institutions must undertake, including broad education. He feared lifelong repetitive labor would dull people, with education counteracting the damage.

These additions reveal a fuller Smith: self-interest is a powerful engine, but it needs the chassis of \textbf{institutions and moral sentiments}. Sympathy makes others' judgments matter; law blocks monopoly and fraud; education supports people worn down by specialization. An engine without this chassis was not his desired society. Reducing Smith to selfishness-is-right wrongs the author and demonstrates why primary reading matters: \textbf{returning to a text often overturns popular paraphrase}.

\section{One Factory, Three Explanations}
Return to Manchester with our components and primary texts. We can compare explanations of one fact: fourteen-hour assembly-line labor earns subsistence wages while owners receive profits. Different frameworks tell different stories. Keep three layers distinct: hours, wages, and profit figures are verifiable evidence; why this happens and what it means are disputed explanation; how to judge it is value, best considered after explanation. We compare the second layer here.

\textbf{Explanation one: division of labor and exchange, in Smith's tradition.}

Wealth comes from specialization: ten people each handling a stage make more tea than one working alone. But somebody must advance funds. Workers cannot wait for drinks to sell, so capitalists supply equipment, materials, and wages in advance, enabling specialization. Profit rewards funding, organization, and risk; wages price labor time through markets. Mary earns little because replacements are plentiful; Robert earns much because few dare make his investment.

Its strength is explaining wealth's \textbf{creation} and factories' genuinely cheaper, more widely affordable cloth. Its limit is seeing supply-and-demand pricing everywhere while overlooking power behind prices: when Mary accepts low wages, refusal is not actually an option. An explanation missing lack of choice misses half of reality.

\textbf{Explanation two: exploitation, in Marx's tradition.}

Here labor creates cloth's value; machinery and cotton merely transfer existing value into the product. Daily labor creates more value than daily wages. Capital appropriates the difference, Marx's \textbf{surplus value}, without compensation. Profit is not payment for advances but labor's creation assigned institutionally to capital. Marx and Engels famously described this in the 1848 Communist Manifesto: the bourgeoisie reduced human relationships to naked interest and merciless cash transactions.

Its strength is focusing for the first time on profit's origin, passed lightly over by the first account, and explaining why those creating most often receive least: institutions leave them only labor time to sell. Its limits are real too: attributing all value to labor struggles to accommodate scarce risk bearing, organization, and technological judgment. Some predictions also imply capitalism should long since have collapsed, yet it repeatedly transforms and survives, as the next section examines.

\textbf{Explanation three: institutional construction, in Polanyi's tradition.}

Twentieth-century Hungarian-born Karl Polanyi offered a third story in The Great Transformation. Labor, land, and money were not originally made for sale: people are living beings, land nature and home, money a bookkeeping tool. Turning them into commodities was not spontaneous but \textbf{constructed} through political decisions. Enclosures drove English peasants from land; the 1834 New Poor Law pushed the poor into mills. The free market did not grow unaided: the state chiseled it into existence.

Its strength is explaining no choice: Mary's predicament follows changeable political arrangements, not cold supply-and-demand laws. What was made can be remade. Its limit is emphasizing market society's construction while underplaying exchange's own expansive logic. Modern politicians did not invent human exchange; Sima Qian saw it two millennia earlier.

These explanations are not simply mutually right or wrong. One explains creation, another unequal distribution, the third rules' origins. A mature thinker carries more than one hammer for large problems.

\section[What Lives Inside a Commodity?]{Deeper Challenge: What Lives Inside a Commodity?}
Having mentioned Marx three times, formally introduce his theory, one of the humanities and social sciences' weightiest analytical tools, worth the effort. It lies between economics and sociology. His principal work is Capital, whose first volume appeared in 1867.

It begins with an unassuming \textbf{commodity}: milk tea, sneakers, game skins. Marx gives every commodity two faces. \textbf{Use value} means usefulness: tea quenches thirst and shoes protect feet. \textbf{Exchange value} means what it trades for, its monetary worth. These diverge: water's use may mean life or death yet cost two yuan, while limited-edition shoes may offer ordinary utility at dozens of times the price. Price never simply says how important something is to you.

What determines exchange value? Marx pushes classical economics to its conclusion: ultimately, \textbf{socially necessary labor time}, the labor required under average social technology and skill. Machinery and materials transfer their value; only labor creates additional value.

The crucial next move is that \textbf{labor power itself becomes a commodity}. Strictly, Mary sells not labor but a day's \textbf{labor power}, her capacity to work. This commodity uniquely creates more value in use than its own value. Maintaining her daily capacity---food, clothing, housing, family---costs the equivalent of a wage; her day's cloth sells for more. Marx calls the difference surplus value, appearing as profit, interest, and rent.

This yields \textbf{capital accumulation}. Capital's full journey is money $\rightarrow$ buying materials and labor power $\rightarrow$ producing commodities $\rightarrow$ selling $\rightarrow$ more money. Reinvest the increment and the cycle expands. Accumulation is capital's nature: stopped capital ceases to be capital. Growth becomes compulsory, not optional; expanding firms swallow nonexpanding ones.

Fourth comes \textbf{class}. Marx does not simply mean rich and poor, descriptions of wallet thickness. Class concerns your \textbf{position} in production: living by owning means of production---factories, equipment, capital---or selling labor power. Frugal and extravagant owners share a class, as do optimistic and angry Marys. Class analysis locates people, not their moral character.

The subtlest final concept is \textbf{commodity fetishism}. The tea's eighteen-yuan label shows cup, pearls, and milk foam, but hides the worker's back strain, landlord's bargaining power, and franchise contract. Human relationships---employment, renting, commission---appear as \textbf{relationships between things}: tea worth eighteen yuan. Things' clothing conceals people's circumstances. Discussing primitive accumulation through enclosure, colonialism, and slave trading in Capital's first volume, Marx wrote a severe sentence, rendered in standard Chinese translations as:

\begin{quote}
Capital comes into the world dripping blood and filth from head to foot, from every pore.
\end{quote}
Remember that this concerns the \textbf{starting point}: much initial wealth came from enclosure, colonial plunder, and slavery. Read it beside Manchester's dawn for a complete chain: wealth accumulates, costs are distributed, and both disappear behind price tags.

The theory remains used, disputed, and revised. It cannot explain everything, as its limits show, but it gave the world eyes it cannot remove: see a commodity and ask about the people behind it.

One last piece: \textbf{capitalism has never been a fixed machine}. Sixteenth- and seventeenth-century Venice and Amsterdam practiced \textbf{commercial capitalism}, profiting chiefly from trade and finance. Founded in 1602, the Dutch East India Company even sold public shares, letting ordinary people buy a piece of an overseas trading empire for the first time. Nineteenth-century Manchester exemplifies \textbf{industrial capitalism}, hiring labor for machine production. Banks and corporations joined in late-nineteenth-century \textbf{finance capitalism}, making money grow faster than factories. After depression and two world wars, mid-twentieth-century Western states added welfare, labor law, and public education, with Nordic countries going furthest. East Asian governments deeply planned industry, producing state-led forms. Today's platforms, taking commissions on content they do not produce and coordinating millions of drivers without owning vehicles, are often called \textbf{platform capitalism}. Components remain, assemblies change. History first refutes anyone praising or cursing an unchanging statue: capitalism resembles an organism continually changing costume.

\section[You Are Inside This Milk Tea]{Back to Today: You Are Already Inside This Milk Tea}
This nineteenth-century question sits in your pocket now.

Delivery riders' time is divided algorithmically into minute-sized boxes: assignments, countdowns, lateness fines. Manchester's whistle moved into phones, ubiquitous enough to need no closing gate. Jokes about 996 schedules and wage workers resonate because countless people recognize Mary: sellers of time have limited autonomy over it.

Consider consumption too. Sneakers, blind boxes, game skins, and character merchandise reveal diverging use and exchange values. How much of a shoe's price pays for wearing it and how much for its symbolism? More subtly, when you watch free short videos, you are the commodity rather than customer: platforms package your attention for advertisers. You now have enough components to unpack commodity fetishism yourself.

Schools echo it. Involution names an accumulation logic: everyone spends longer drilling questions, admission cutoffs rise, yet relative positions scarcely change. An arms race of effort projects expand-or-be-swallowed accumulation into education.

But give the other side its full case or learn only half the tool. This mechanism reliably supplies eighteen-yuan teas in mall basements, lets an unemployed person become a rider the next day, and turns grandparents' rationed cloth into today's overwhelming clothing choice. Efficiency and costs are both real, as Manchester showed. Hear one more voice through: eight-hour days, weekends, child-labor bans, and compulsory schooling were won through arguments, strikes, and laws. Manchester's dawn was not the endpoint. Human action changes institutions. Polanyi was right: what was made can be remade.

\section{Your Question: Should This Milk Tea Exist?}
Return to the eighteen-yuan tea and close the accounts yourself.

Suppose the owner's profit is abolished and all tea revenue goes to workers. Fair, perhaps---but who funds the next shop, replaces broken equipment, or covers slow-season losses? If workers contribute, should the largest contributor receive more? How does that additional share differ from the owner's former profit?

Conversely, keep profit for funders and hourly time sales for workers. Can workers accumulate capital and become funders? If this road is nearly closed to most, how much weight does voluntary agreement retain?

Your task is to \textbf{design a tea-shop distribution rule that makes operation possible and seems fair to you}. Explain who funds it, bears risk, receives profit, and has a voice, and whether your rule suits one shop or the whole neighborhood.

Give the rule and reasons. There is no standard answer. Every should you write answers a larger question: what entitles a person to something---labor contributed, risk borne, property owned, or human dignity? Asked from Sima Qian's market to Manchester's mill and your tea shop, that question remains unfinished.
